\subsection{Introducing Inequalities} 
We have several symbols that we use to describe relationships between numbers:
\begin{center}
\begin{tabular}{cl}
$=$&is equal to (or equals)\\
$\neq$&is not equal to\\
$<$&is less than\\
$\leq$&is less than or equal to\\
$>$&is greater than\\
$\geq$&is greater than or equal to
\end{tabular}
\end{center}
Some notes about these symbols:
\begin{itemize}
\item We understand ``is less than'' to mean ``is to the \makeblank[1in] on the number line,'' and ``is greater than'' to mean ``is to the \makeblank[1in] on the number line.'' So $3\makeblanksmall 5$, but $-3\makeblanksmall -5$.
\item For fixed numbers, exactly one of \makeblanksmall, \makeblanksmall, or \makeblanksmall will always apply.
\item When we write an inequality with a \makeblank[1.5in], we are describing many numbers.
\end{itemize}
\begin{definition}[Solution of an Inequality]
A \term{solution} of an equation is a value (or values) of the variable (or variables) that make the \makeblank[1.5in] true.
\end{definition}
\begin{Example}
List some numbers that are solutions to the inequality $x<5$.
\begin{cpic}{}{2}
\end{cpic}
\end{Example}
\begin{Example}
List some numbers that are solutions to the inequality $x\leq 5$. What is the difference between $x<5$ and $x\leq 5$?
\begin{cpic}{}{2}
\end{cpic}
The number \makeblanksmall is a solution to $x\leq 5$ and not $x<5$.
\end{Example}
We can use inequalities to describe numbers being positive or negative.
\begin{itemize}
\item To say that $x$ is positive, we write \makeblank[1in].
\item To say that $x$ is nonnegative (positive or zero), we write \makeblank[1in].
\item To say that $x$ is negative, we write \makeblank[1in].
\item To say that $x$ is nonpositive (negative or zero), we write \makeblank[1in].
\end{itemize}